\documentclass[11pt]{article}

\usepackage{series}

\renewcommand\SeriesName{Modal Triplet Theory}
\renewcommand\SeriesShort{Admissible Projection}
\renewcommand\PartRoman{I}
\renewcommand\PartArabic{1}
\renewcommand\PartsInSeries{I}

\title{\Large Dirac Delta Functions as Singular Shadows of Admissible Projection\\
\large A Fixed-Point and Gauge-Theoretic Formulation in Modal Triplet Theory}
\author{Peter Nero}
\date{April 2026}

\hypersetup{
  pdftitle={Dirac Delta Functions as Singular Shadows of Admissible Projection},
  pdfauthor={Peter Nero},
  pdfsubject={Modal Triplet Theory; admissible projection; Dirac delta; fixed points; gauge fixing},
  pdfkeywords={Modal Triplet Theory, Dirac delta, spectral projector, heat kernel, gauge fixing, Faddeev-Popov, fixed points}
}

\begin{document}
\maketitle
\seriestagline

\begin{abstract}
Dirac delta distributions appear throughout physics as identity kernels, point sources,
constraint enforcers, measurement idealizations, gauge-fixing devices, and conservation
laws at interaction vertices. Their ubiquity is usually treated as a technical feature of
continuum mathematics. This paper isolates a structural explanation compatible with
Modal Triplet Theory (MTT): a Dirac delta is the singular downstream limit of an
admissible projection when finite-width selection, coherent-sector truncation, or
representative choice is idealized to zero width.

The mathematical core is deliberately narrow. We prove that spectral projection kernels
on a compact elliptic setting converge to the Dirac delta in the distributional sense, and
we record the corresponding heat-kernel approximate identity. These standard analytic
facts provide the rigorous anchor for the MTT interpretation: the native object in a
finite-capacity coherent regime is not the exact identity kernel \(\delta(x-y)\), but a
bounded coherent kernel \(K_{\mathrm{coh}}(x,y)\) induced by a coherent projector
\(\Pi_{\mathrm{coh}}\). The delta is recovered only after an infinite-resolution or
zero-width limit.

We then apply this projection reading to Green kernels, canonical commutators, contact
interactions, path-integral constraints, measurement, fixed-point basin selection, and
gauge fixing. Gauge theory displays the same architecture explicitly: gauge freedom is
uncollapsed projection redundancy, gauge fixing is local representative selection, the
gauge-fixing delta is a singular section-selection kernel, and the Faddeev--Popov
determinant is the projection Jacobian. The paper does not replace distribution theory,
QFT, or gauge theory. It supplies a diagnostic principle: wherever a physical theory writes
a Dirac delta, one should ask which finite admissible projection, filter, basin, or quotient
operation has been idealized as exact.
\end{abstract}

\tableofcontents

\section{Purpose and claim discipline}

The purpose of this paper is to isolate a recurring mathematical pattern in physical
theories and reinterpret it through the projection-first architecture of MTT.

The recurring pattern is the Dirac delta. It appears as:
\begin{enumerate}
\item the identity kernel \(\delta(x-y)\);
\item a point source \(\delta(x-x_0)\);
\item a hard constraint \(\delta(C[x])\);
\item a sharp measurement outcome;
\item a gauge-fixing slice selector;
\item a conservation law at interaction vertices;
\item a spectral selector \(\delta(E-E_n)\);
\item a contact interaction or coincidence-limit singularity.
\end{enumerate}

The central claim is structural:
\begin{center}
\fbox{\begin{minipage}{0.86\linewidth}
A Dirac delta marks where a finite admissible process has been idealized as exact
zero-width selection.
\end{minipage}}
\end{center}

\subsection{Non-claims}

This paper does not claim:
\begin{enumerate}
\item that distribution theory is mathematically invalid;
\item that ordinary QFT must be reformulated from scratch;
\item that the coherent-kernel replacement is numerically evaluated here for all physical sectors;
\item that every delta function has the same microscopic realization;
\item that the proposed reading alone solves renormalization, measurement, or quantum gravity.
\end{enumerate}

The narrower claim is:

\begin{quote}
In any MTT-compatible effective description, a Dirac delta should be read as a
singular encoding of a bounded projection, admissibility filter, survivor-basin
selection, or representative-choice operation whenever that effective description
arises by coherent truncation or quotienting.
\end{quote}

\subsection{Working dictionary}

The following dictionary states the intended use of the paper before the analytic details are
introduced.  The left column records a standard distributional device; the right column gives
the corresponding projection-theoretic reading.

\begin{center}
\begin{tabular}{p{0.36\linewidth}p{0.52\linewidth}}
\hline
\textbf{Standard delta usage} & \textbf{Projection/admissibility reading} \\
\hline
\(\delta(x-y)\) & Exact identity kernel; singular limit of a coherent identity kernel.\\
\(\delta(x-x_0)\) & Point localization; zero-width limit of a finite survivor basin or source profile.\\
\(\delta(C[\phi])\) & Hard constraint; zero-width limit of an admissibility filter.\\
\(\delta(G[A])\) & Gauge-slice selector; singular representative choice in a quotient.\\
\(\delta(\sum_i p_i)\) & Exact bookkeeping closure at a vertex; sharp limit of finite overlap conservation.\\
\(|x\rangle\langle x|\) & Ideal measurement effect; zero-width limit of finite detector/coherence support.\\
\hline
\end{tabular}
\end{center}

This table is not a proof.  It is a map of the examples to which the proved kernel theorem
will later be applied.


\section{Fixed-point backbone: projection before delta}

The fixed-point framework supplies the analytic backbone for the present proposal.  The
basic structure is
\begin{equation}
T_t=\Pi_{\mathrm{coh}}\circ \Phi_t,
\end{equation}
where \(\Phi_t\) is a smoothing or dissipative flow on the underlying modal configuration
space and \(\Pi_{\mathrm{coh}}\) is a bounded coherent projector onto the joint harmonic or
coherent sector.  The map \(T_t\) is the projected evolution whose fixed points define
stable coherent regimes.

If \(\Pi_{\mathrm{coh}}\) admits an integral kernel \(K_{\mathrm{coh}}\), then
\begin{equation}
(\Pi_{\mathrm{coh}}f)(x)=\int_X K_{\mathrm{coh}}(x,y)f(y)\,\dd y.
\end{equation}
Standard continuum physics often uses instead
\begin{equation}
f(x)=\int_X \delta(x-y)f(y)\,\dd y.
\end{equation}
The difference is the difference between a finite admissible identity and an ideal exact
identity.  The MTT replacement principle is therefore
\begin{equation}
\boxed{\delta(x-y)\quad\leadsto\quad K_{\mathrm{coh}}(x,y).}
\end{equation}

\begin{definition}[Finite coherent identity kernel]
Let \(\Pi_{\mathrm{coh}}\) be a bounded projection on a Hilbert space of fields over a
domain \(X\).  If \(\Pi_{\mathrm{coh}}\) is represented by a distributional or smooth kernel
\(K_{\mathrm{coh}}(x,y)\), then \(K_{\mathrm{coh}}\) is called the finite coherent identity
kernel of that effective regime. It acts as an identity only on \(\Ran(\Pi_{\mathrm{coh}})\).
\end{definition}

\begin{remark}
The phrase ``identity kernel'' is sector-relative.  A finite coherent kernel need not be
the identity on the full upstream function space.  It is the identity on the retained
admissible sector and a filter on everything else.
\end{remark}

\section{Spectral gaps and finite width}

The fixed-point framework assumes a positive spectral separation between coherent and
incoherent sectors.  In a standard elliptic model,
\begin{equation}
A|_{\Ran(Q)}\geq \lambda^\ast>0,\qquad Q=\Id-\Pi_{\mathrm{coh}}.
\end{equation}
This yields damping estimates of the form
\begin{equation}
\norm{A^{1/2}e^{-tA}Q}\lesssim t^{-1/2}e^{-\lambda^\ast t}.
\end{equation}

The interpretation used below is:
\[
\boxed{\text{finite spectral gap and finite damping margin imply finite projection width.}}
\]
The delta limit corresponds to an idealization in which coherent projection becomes
infinitely sharp.  MTT does not begin with this sharpness.  It begins with bounded
projection, finite spectral separation, and controlled truncation.

A canonical model for the replacement is the heat kernel
\begin{equation}
H_\tau(x,y)\sim (4\pi \tau)^{-d/2}\exp\left(-\frac{\dist(x,y)^2}{4\tau}\right),
\end{equation}
with
\begin{equation}
\delta(x-y)=\lim_{\tau\downarrow 0}H_\tau(x,y)
\end{equation}
in the distributional sense.

\section{Spectral projection kernels and the Dirac delta}

This section supplies the narrow mathematical anchor.  The aim is not to prove that every
occurrence of a Dirac delta in physics has the same origin.  The aim is to prove the
precise analytic statement that a Dirac delta arises as the distributional limit of
increasingly complete projection kernels.

\begin{assumption}[Compact elliptic setting]\label{ass:compact-elliptic}
Let \((X,g)\) be a compact smooth Riemannian manifold without boundary, or a compact
domain with boundary equipped with a self-adjoint elliptic boundary condition. Let
\(\Delta\geq 0\) be a nonnegative Laplace-type operator on \(L^2(X)\).  Let
\[
0\leq \lambda_0\leq \lambda_1\leq \lambda_2\leq \cdots
\]
be its eigenvalues, repeated with multiplicity, and let \(\{\phi_n\}_{n=0}^\infty\) be an
orthonormal eigenbasis satisfying
\begin{equation}
\Delta \phi_n=\lambda_n\phi_n .
\end{equation}
\end{assumption}

For \(\Lambda>0\), define the spectral projector
\begin{equation}
\Pi_\Lambda f=\sum_{\lambda_n\leq \Lambda}\ip{\phi_n}{f}\phi_n .
\end{equation}
Its Schwartz kernel is
\begin{equation}
K_\Lambda(x,y)=\sum_{\lambda_n\leq \Lambda}\phi_n(x)\overline{\phi_n(y)} .
\end{equation}
Then
\begin{equation}
(\Pi_\Lambda f)(x)=\int_X K_\Lambda(x,y)f(y)\,\dd\mathrm{vol}_g(y).
\end{equation}

\begin{theorem}[Spectral projection kernels converge to the Dirac delta]\label{thm:spectral-delta}
Under \cref{ass:compact-elliptic}, for every \(f\in C^\infty(X)\),
\begin{equation}
\lim_{\Lambda\to\infty}\int_X K_\Lambda(x,y)f(y)\,\dd\mathrm{vol}_g(y)=f(x),
\end{equation}
with convergence in \(C^\infty(X)\).  Equivalently,
\begin{equation}
K_\Lambda(x,y)\longrightarrow \delta(x-y)
\end{equation}
in the sense of distributions on \(X\times X\), where \(\delta(x-y)\) denotes the kernel
of the identity operator with respect to the Riemannian volume measure.
\end{theorem}

\begin{proof}
By the spectral theorem, \(\{\phi_n\}\) is a complete orthonormal basis of \(L^2(X)\).
Hence every \(f\in L^2(X)\) has the expansion
\[
f=\sum_{n=0}^{\infty} f_n\phi_n,\qquad f_n=\ip{\phi_n}{f},
\]
with convergence in \(L^2\).  For \(f\in C^\infty(X)\), elliptic regularity gives rapid
decay of spectral coefficients: for every integer \(m\geq 0\),
\begin{equation}
\sum_{n=0}^{\infty}(1+\lambda_n)^m|f_n|^2<\infty .
\end{equation}

The truncated projection is
\[
\Pi_\Lambda f=\sum_{\lambda_n\leq \Lambda}f_n\phi_n,
\]
so
\[
f-\Pi_\Lambda f=\sum_{\lambda_n>\Lambda}f_n\phi_n.
\]
For every Sobolev index \(s\geq 0\),
\begin{equation}
\norm{f-\Pi_\Lambda f}_{H^s}^2
=\sum_{\lambda_n>\Lambda}(1+\lambda_n)^s|f_n|^2 .
\end{equation}
The right-hand side tends to zero as \(\Lambda\to\infty\), because \(f\) is smooth and
therefore has finite \(H^s\)-norm for every \(s\).  Hence \(\Pi_\Lambda f\to f\) in all
Sobolev norms.  By Sobolev embedding, the convergence is in \(C^k\) for every finite
\(k\).  Therefore
\[
\lim_{\Lambda\to\infty}(\Pi_\Lambda f)(x)=f(x)
\]
smoothly in \(x\).

Since
\[
(\Pi_\Lambda f)(x)=\int_X K_\Lambda(x,y)f(y)\,\dd\mathrm{vol}_g(y),
\]
we obtain the stated convergence.  This is precisely the distributional defining property
of the Dirac delta kernel.  Therefore \(K_\Lambda\to\delta\) distributionally on
\(X\times X\).
\end{proof}

\begin{corollary}[Finite projection kernels are regularized deltas]\label{cor:finite-proj}
For finite \(\Lambda\), \(K_\Lambda(x,y)\) is a finite-rank smooth kernel. It acts as the
identity on \(\Ran(\Pi_\Lambda)\) and as a truncation on the full space.  Thus
\(K_\Lambda\) is a finite-resolution identity kernel, while \(\delta(x-y)\) is the
infinite-bandwidth identity kernel.
\end{corollary}

\begin{proof}
Finite-rank smoothness follows from the finite sum defining \(K_\Lambda\).  The identity
property on \(\Ran(\Pi_\Lambda)\) follows from idempotence:
\(\Pi_\Lambda^2=\Pi_\Lambda\).  The limiting statement is \cref{thm:spectral-delta}.
\end{proof}

\begin{corollary}[MTT reading]\label{cor:mtt-reading}
If the coherent sector of an MTT fixed-point regime is represented by a bounded spectral
projector \(\Pi_{\mathrm{coh}}\), then its kernel
\begin{equation}
K_{\mathrm{coh}}(x,y)=\langle x|\Pi_{\mathrm{coh}}|y\rangle
\end{equation}
is the MTT-native replacement for the exact identity kernel \(\delta(x-y)\) on that
regime.  The Dirac delta is recovered only when the coherent projection is idealized as
complete, infinitely sharp, and free of finite-capacity remainder.
\end{corollary}

\begin{remark}
\Cref{cor:mtt-reading} is interpretive rather than an additional analytic theorem.  The
analytic theorem is the distributional convergence of spectral projection kernels.  The MTT
claim is that physical deltas should be read through this projection structure whenever the
relevant effective theory arises from coherent truncation.
\end{remark}

\begin{center}
\fbox{\begin{minipage}{0.88\linewidth}
\textbf{Scope of proof.}  The result proved in this section is the standard analytic fact
that increasing spectral projection kernels converge to the identity kernel, i.e. to the
Dirac delta distribution.  The paper does \emph{not} prove here that every physical delta
is generated by MTT projection.  The broader claim is a diagnostic and interpretive one:
when an effective physical description is known or assumed to arise from coherent
truncation, quotienting, or admissible selection, its delta functions should be examined as
singular limits of the corresponding finite kernels.
\end{minipage}}
\end{center}

\section{Heat kernels as admissibility kernels}

Spectral cutoffs are not the only regularized identity kernels.  Heat kernels provide the
smoothest and most physically useful model.

Let \(H_\tau(x,y)\) be the heat kernel of \(\Delta\):
\begin{equation}
e^{-\tau\Delta}f(x)=\int_X H_\tau(x,y)f(y)\,\dd\mathrm{vol}_g(y).
\end{equation}
Spectrally,
\begin{equation}
H_\tau(x,y)=\sum_{n=0}^{\infty}e^{-\tau\lambda_n}\phi_n(x)\overline{\phi_n(y)}.
\end{equation}

\begin{theorem}[Heat-kernel approximate identity]\label{thm:heat}
For every \(f\in C^\infty(X)\),
\begin{equation}
\lim_{\tau\downarrow 0}\int_X H_\tau(x,y)f(y)\,\dd\mathrm{vol}_g(y)=f(x)
\end{equation}
in \(C^\infty(X)\).  Equivalently,
\begin{equation}
H_\tau(x,y)\to \delta(x-y)
\end{equation}
distributionally as \(\tau\downarrow 0\).
\end{theorem}

\begin{proof}
The heat semigroup \(e^{-\tau\Delta}\) converges strongly to the identity on every Sobolev
space \(H^s(X)\).  For smooth \(f\), this convergence holds in all Sobolev norms and hence,
by Sobolev embedding, in \(C^k\) for every \(k\).  The kernel formulation gives the stated
distributional convergence.
\end{proof}

\begin{remark}[Projection versus heat smoothing]
A sharp spectral projector \(\Pi_\Lambda\) gives a band-limited coherent identity.  A heat
kernel \(H_\tau\) gives a soft coherent identity with exponential high-frequency suppression.
Both converge to \(\delta\) in singular limits.  In MTT applications, heat/proper-time kernels
are often more natural because finite damping and finite coherence capacity are represented
by exponential suppression rather than a hard spectral wall.
\end{remark}

\section{Inverse principle: delta as diagnostic}

\begin{definition}[Delta diagnostic]\label{def:delta-diagnostic}
The delta diagnostic is the following rule: whenever a physical theory contains a Dirac
delta, ask which finite projection, admissibility filter, survivor-basin selection,
representative choice, or coherent identity kernel has been idealized to zero width.
\end{definition}

This principle does not deny standard calculations.  It identifies where the calculation has
hidden a projection step.

\begin{example}[Green kernel]
The equation
\begin{equation}
LG(x,y)=\delta(x-y)
\end{equation}
states that \(G\) is the response to a perfectly localized source.  The coherent replacement is
\begin{equation}
LG_{\mathrm{coh}}(x,y)=K_{\mathrm{coh}}(x,y).
\end{equation}
Thus the point source is replaced by an admissibly localized source.
\end{example}

\begin{example}[Canonical commutator]
The standard equal-time relation
\begin{equation}
[\phi(t,x),\pi(t,y)]=i\hbar\delta(x-y)
\end{equation}
is replaced, in a projected coherent sector, by
\begin{equation}
[\phi_{\mathrm{coh}}(t,x),\pi_{\mathrm{coh}}(t,y)]
=i\hbar K_{\mathrm{coh}}(x,y).
\end{equation}
Locality is not destroyed.  It becomes coherent-sector-local rather than infinitely sharp.
\end{example}

\section{Fixed points and delta concentration}

In ordinary dynamical systems, convergence to a stable fixed point is often represented
distributionally as
\begin{equation}
\rho_t\to \delta_{x_\ast}.
\end{equation}
In MTT this should be refined.  A fixed point \(\Psi_\ast\) of
\begin{equation}
T_t=\Pi_{\mathrm{coh}}\circ\Phi_t
\end{equation}
is the center of a stabilized coherent basin, not necessarily a literal zero-width state of the
full upstream dynamics.

The MTT-native statement is
\begin{equation}
\rho_t\to \rho_{\mathrm{basin},\Psi_\ast},
\end{equation}
where \(\rho_{\mathrm{basin},\Psi_\ast}\) has finite width determined by projection resolution,
damping margins, disturbance floors, and the size of the basin of attraction.  The delta
appears only in the singular limit:
\begin{equation}
\rho_{\mathrm{basin},\Psi_\ast}\to\delta_{\Psi_\ast}.
\end{equation}

\begin{proposition}[Projected fixed-point delta]\label{prop:fixed-delta}
Suppose an admissible projected dynamics \(T_t\) has an attracting fixed point \(\Psi_\ast\)
with a basin whose effective distribution at time \(t\) is \(\rho_t\). If \(\rho_t\) converges
weakly to a probability measure \(\rho_\ast\) supported in a finite coherent basin
\(B_\ast\), then the notation \(\rho_t\to\delta_{\Psi_\ast}\) is valid only after the further
idealization that \(B_\ast\) has zero effective width.
\end{proposition}

\begin{proof}
Weak convergence to \(\delta_{\Psi_\ast}\) means that all continuous test observables take the
value associated with the single point \(\Psi_\ast\).  If the limiting measure is supported on
a finite basin \(B_\ast\), then observables that vary across \(B_\ast\) distinguish
\(\rho_\ast\) from \(\delta_{\Psi_\ast}\).  The delta notation is therefore justified only when
the retained observable algebra cannot resolve the basin width, or in the additional
zero-width limit.
\end{proof}

\section{Disturbance--damping balance and nonzero width}

The fixed-point disturbance analysis supports the finite-width replacement.  For a
non-harmonic mode with damping margin \(\gamma_{n,k}>0\) and disturbance strength
\(\delta_{n,k}\), the Ornstein--Uhlenbeck regime gives
\begin{equation}
\sigma^2_{n,k}=\frac{\delta_{n,k}}{2\gamma_{n,k}}.
\end{equation}
Thus exact delta collapse requires either
\begin{equation}
\delta_{n,k}\to 0
\end{equation}
or
\begin{equation}
\gamma_{n,k}\to\infty.
\end{equation}
Neither is generic in finite-capacity MTT.  A Gaussian model for the finite stabilized kernel is
\begin{equation}
K_{\gamma,\delta}(x,x_0)\sim
\exp\left[-\frac{(x-x_0)^2}{2\sigma^2}\right],
\qquad
\sigma^2=\frac{\delta}{2\gamma}.
\end{equation}

\section{Gauge as uncollapsed projection; delta as collapsed projection}

Gauge theory provides the clearest standard example of the same architecture.

Let \(\mathcal A\) be a space of gauge fields and \(\mathcal G\) the gauge group.  The
physical configuration space is the quotient
\begin{equation}
\mathcal A/\mathcal G.
\end{equation}
The projection
\begin{equation}
\mathcal A\to\mathcal A/\mathcal G
\end{equation}
is many-to-one.  Gauge freedom is the visible persistence of that non-injectivity:
\[
\boxed{\text{gauge freedom}=\text{uncollapsed projection redundancy}.}
\]
In MTT language, this is lens structure.

Gauge fixing imposes a condition
\begin{equation}
G[A]=0.
\end{equation}
This selects a representative slice through each gauge orbit.  The Faddeev--Popov identity
has the formal form
\begin{equation}
1=\Delta_{\mathrm{FP}}[A]\int \mathcal D\alpha\,
\delta(G[A^\alpha]).
\end{equation}
Here \(A^\alpha\) parameterizes a gauge orbit, \(\delta(G[A^\alpha])\) selects the slice,
and \(\Delta_{\mathrm{FP}}\) corrects the quotient measure.

\begin{proposition}[Faddeev--Popov determinant as projection Jacobian]\label{prop:fp}
Assume a local gauge slice \(G[A]=0\) intersects gauge orbits transversely near \(A\).  Then
the Faddeev--Popov factor
\begin{equation}
\Delta_{\mathrm{FP}}[A]
=\left|\det D_\alpha(G[A^\alpha])\right|
\end{equation}
is the Jacobian of the projection from orbit coordinates to the chosen local slice.
\end{proposition}

\begin{proof}
Locally decompose a neighborhood of \(A\) into coordinates \((s,\alpha)\), where \(s\)
parametrizes the gauge slice and \(\alpha\) parametrizes the gauge orbit.  The gauge condition
maps orbit coordinates to the constraint value \(G[A^\alpha]\).  If the intersection is
transverse, \(D_\alpha G[A^\alpha]\) is invertible in the orbit directions.  The standard
change-of-variables formula then contributes the absolute determinant
\(\left|\det D_\alpha(G[A^\alpha])\right|\), which is precisely the Faddeev--Popov factor.
\end{proof}

Thus:
\[
\boxed{\delta(G[A])=\text{singular shadow of representative selection},}
\]
and
\[
\boxed{\Delta_{\mathrm{FP}}[A]=\text{projection Jacobian}.}
\]

\subsection{Ghosts as quotient bookkeeping}

The determinant may be represented by ghost fields:
\begin{equation}
\Delta_{\mathrm{FP}}[A]=
\int \mathcal D\bar c\,\mathcal Dc\,e^{iS_{\mathrm{ghost}}[\bar c,c,A]}.
\end{equation}
In MTT terms:
\[
\boxed{\text{ghosts encode the local measure cost of quotienting lens redundancy}.}
\]
They are not physical particles in the ordinary sense.  They are bookkeeping fields for the
projection Jacobian.

\subsection{Gribov copies}

Gauge fixing may fail globally. A gauge slice may intersect one orbit multiple times.  These
are Gribov copies. In MTT language:
\[
\boxed{\text{Gribov copies}=\text{failure of global admissible section}.}
\]
This matches the MTT principle that reduced descriptions are local encodings, not globally
valid ontologies.

\subsection{BRST}

BRST symmetry algebraically controls the quotient:
\begin{equation}
Q_{\mathrm{BRST}}^2=0,\qquad
\mathcal H_{\mathrm{phys}}=\ker Q_{\mathrm{BRST}}/\im Q_{\mathrm{BRST}}.
\end{equation}
MTT reading:
\[
\boxed{\text{physical states}=\text{admissible quotient classes modulo null redundancy}.}
\]

\section{Path integrals and admissibility filters}

Path integrals often impose constraints using hard deltas:
\begin{equation}
Z=\int \mathcal D\phi\,\delta(C[\phi])e^{iS[\phi]/\hbar}.
\end{equation}
MTT replaces this with a finite admissibility kernel:
\begin{equation}
\delta(C[\phi])\quad\leadsto\quad \mathcal K_{\mathrm{adm}}[C[\phi]].
\end{equation}
A standard Gaussian model is
\begin{equation}
\mathcal K_{\mathrm{adm}}[C]
=\exp\left(-\frac{1}{2\epsilon_{\mathrm{adm}}^2}\norm{C}^2\right).
\end{equation}
Thus
\begin{equation}
Z_{\mathrm{MTT}}=
\int \mathcal D\phi\,
\exp\left(\frac{i}{\hbar}S[\phi]
-\frac{1}{2\epsilon_{\mathrm{adm}}^2}\norm{C[\phi]}^2\right).
\end{equation}
The hard delta is recovered as
\begin{equation}
\delta(C)=\lim_{\epsilon_{\mathrm{adm}}\to 0}\mathcal K_{\mathrm{adm}}[C].
\end{equation}

\section{Measurement as finite survivor-basin selection}

An ideal position measurement uses projectors
\begin{equation}
P_x=|x\rangle\langle x|,
\end{equation}
with distributional kernel
\begin{equation}
\langle y|P_x|z\rangle=\delta(y-x)\delta(z-x).
\end{equation}
MTT replaces this with a finite coherent effect:
\begin{equation}
E_x^{\mathrm{coh}}(y,z)
=K_{\mathrm{coh}}(y,x)\overline{K_{\mathrm{coh}}(z,x)}.
\end{equation}
The probability becomes
\begin{equation}
p(x)=\langle \psi|E_x^{\mathrm{coh}}|\psi\rangle.
\end{equation}
The post-measurement state is not a point delta but a stabilized survivor basin:
\begin{equation}
\psi\mapsto \frac{M_x^{\mathrm{coh}}\psi}{\norm{M_x^{\mathrm{coh}}\psi}}.
\end{equation}
Thus:
\[
\boxed{\text{measurement collapse}=\text{finite basin selection idealized as delta collapse}.}
\]

\section{QFT applications}

\subsection{Green functions}

Standard:
\begin{equation}
LG(x,y)=\delta(x-y).
\end{equation}
MTT-coherent:
\begin{equation}
LG_{\mathrm{coh}}(x,y)=K_{\mathrm{coh}}(x,y).
\end{equation}

\subsection{Contact interactions}

A local interaction such as \(\lambda\phi^4(x)\) can be written as a zero-width overlap.  The
coherent replacement is
\begin{equation}
\lambda\int dx_1\cdots dx_4\,
V_{\mathrm{coh}}(x_1,x_2,x_3,x_4)
\phi(x_1)\phi(x_2)\phi(x_3)\phi(x_4),
\end{equation}
where
\begin{equation}
V_{\mathrm{coh}}\sim
\int dx\,K_{\mathrm{coh}}(x,x_1)K_{\mathrm{coh}}(x,x_2)
K_{\mathrm{coh}}(x,x_3)K_{\mathrm{coh}}(x,x_4).
\end{equation}
Point-local interaction is the zero-width limit of coherent overlap.

\subsection{Momentum conservation deltas}

At a standard vertex,
\begin{equation}
(2\pi)^4\delta^{(4)}\left(\sum_i p_i\right)
\end{equation}
enforces exact bookkeeping closure.  MTT reads this as the sharp limit of an
overlap-conservation kernel:
\begin{equation}
\delta^{(4)}\left(\sum_i p_i\right)
\quad\leadsto\quad
K_{\mathrm{book}}\left(\sum_i p_i\right),
\end{equation}
where \(K_{\mathrm{book}}\) has width determined by coherent-sector resolution and finite
interaction support.

\section{Renormalization as repair of over-sharp projection}

Many ultraviolet divergences arise from products or limits of distributions at coincident
points:
\begin{equation}
\delta(x-x),\qquad G(x,x),\qquad \phi^n(x)\text{ at one point}.
\end{equation}
The MTT interpretation is:
\begin{center}
\fbox{\begin{minipage}{0.86\linewidth}
UV divergence often signals that a finite-width projection has been replaced by a
zero-width delta.
\end{minipage}}
\end{center}
Renormalization then becomes, at least in part, the downstream repair mechanism required
after over-sharp projection. This does not eliminate renormalization. It explains why
renormalization appears exactly where the idealized continuum theory compresses admissible
overlap structure into pointlike coincidence.

\section{Triadic placement: circle, lens, nil}

The proto-spinor carrier
\begin{equation}
\Xi=(\Psi,C,L,N)
\end{equation}
provides a useful classification.

\begin{center}
\begin{tabular}{ll}
Carrier role & Delta/gauge interpretation\\
\hline
\(C\): circle bookkeeping & phase, return, conservation kernels\\
\(L\): lens redundancy & gauge freedom, equivalence classes, quotienting\\
\(N\): nil survivorship & selection, thresholds, collapse, discrete outcomes
\end{tabular}
\end{center}

Thus:
\begin{equation}
\text{gauge}\subset L,\qquad
\delta\text{-selection}\subset N,\qquad
\text{conservation deltas}\subset C.
\end{equation}
More carefully, delta functions can appear as singular shadows in all three sectors:
circle deltas enforce exact return/bookkeeping closure; lens deltas enforce representative
selection in a redundancy class; nil deltas enforce survivor selection at a threshold.

\section{Diagnostics and possible finite-width effects}

The preceding sections do not by themselves compute new numerical predictions.  They
identify where finite-width corrections would enter if a delta idealization is replaced by an
admissible kernel.  The following are the most concrete diagnostic directions.

\begin{enumerate}
\item \textbf{UV-softened contact interactions.}
A point interaction or local monomial such as \(\lambda\phi^4(x)\) can be replaced by a
finite overlap vertex.  The leading diagnostic is suppression of coincident-point divergences
or a controlled modification of high-momentum behavior.

\item \textbf{Coherent Green functions.}
The equation \(LG=\delta\) is replaced by \(LG_{\mathrm{coh}}=K_{\mathrm{coh}}\).  This gives a
direct worked sequel: compute how a finite coherent source changes the near-source behavior
of a Green function while preserving the ordinary solution outside the kernel width.

\item \textbf{Finite detector-resolution collapse.}
Ideal projectors \(|x\rangle\langle x|\) are replaced by positive finite effects
\(E_x^{\mathrm{coh}}\).  The diagnostic is a nonzero stabilization width rather than exact
delta collapse, with the width controlled by damping and disturbance scales.

\item \textbf{Gauge-fixing tubes instead of exact slices.}
The hard factor \(\delta(G[A])\) can be softened to a tube around the gauge slice.  The
diagnostic is sensitivity to the width of representative selection, especially near Gribov
horizons where the projection Jacobian degenerates.

\item \textbf{Spectral peaks as finite basins.}
A spectral line written as \(\delta(E-E_n)\) is the infinite-lifetime limit of a finite-width
stable mode.  MTT suggests reading linewidths as basin-width or disturbance--damping data
rather than merely as external broadening.
\end{enumerate}

These diagnostics provide the transition from the present foundation paper to concrete
worked examples.  The first technically clean target is the coherent Green-function problem:
replace a point source by a bounded projection kernel and compare the resulting solution to
the ordinary distributional Green function.

\section{Research program}

The practical method is:
\begin{enumerate}
\item Locate every Dirac delta in a physical formulation.
\item Classify its role: identity, source, constraint, representative selection, conservation,
measurement, or spectral selection.
\item Identify the corresponding MTT structure: coherent kernel, admissibility filter, basin
kernel, lens quotient, circle bookkeeping, or nil selection.
\item Replace the delta with a bounded kernel.
\item Study the corrections induced by finite width.
\end{enumerate}

The highest-priority targets are:
\begin{enumerate}
\item canonical commutators;
\item propagators;
\item gauge fixing and Faddeev--Popov determinants;
\item measurement projectors;
\item path-integral constraints;
\item point-particle sources in GR and QFT;
\item contact interactions and UV divergences;
\item spectral delta peaks and finite-lifetime resonances.
\end{enumerate}

\section{Conclusion}

The Dirac delta is one of the most ubiquitous mathematical objects in physics.  In standard
usage it implements identity, localization, constraint, conservation, or selection.  In MTT
these roles share a common structural origin: each is a downstream idealization of projection
under finite admissibility.

The fixed-point framework supplies the analytic core.  Coherent projection is bounded;
smoothing suppresses incoherent modes; spectral gaps control resolution; disturbance--damping
balance leaves finite width; projected fixed points represent stabilized coherent basins.  The
literal Dirac delta appears only when this finite structure is idealized to zero width.

Gauge theory confirms the same pattern from another direction.  Gauge freedom records
uncollapsed projection redundancy; gauge fixing imposes representative selection; the
Faddeev--Popov determinant is the projection Jacobian; ghosts encode quotient-measure
bookkeeping; Gribov ambiguities mark global section failure.

The central principle is:
\[
\boxed{\text{Every Dirac delta is a diagnostic of hidden projection.}}
\]
Where standard physics writes \(\delta\), MTT asks what bounded coherent kernel,
admissibility filter, survivor basin, or gauge-section selection has been collapsed into
singular notation.

\end{document}
